\documentclass{MMCStyle}
\graphicspath{{../figures/}}

\bianhao{202601139}
\tihao{A题}
\timu{长江十年禁渔下鱼类食物链恢复、混沌阈值与入侵风险的耦合建模}
\keyword{长江十年禁渔；四大家鱼；长江江豚；长江鲟；Hastings-Powell模型；李雅普诺夫指数}
\bibliographystyle{gbt7714-numerical}

\begin{document}
\begin{abstract}
围绕长江十年禁渔背景下鱼类生态系统的恢复与失衡问题，本文构建了“资源--功能群模块、珍稀物种模块和三级食物链混沌模块”三层框架。首先，用四类底层资源与四大家鱼功能群的耦合常微分方程刻画禁渔后资源与鱼群的传导；其次，用食源约束、人工放流和通道阻隔共同驱动的模型分析长江江豚与长江鲟的恢复过程；再次，用 Hastings-Powell 三级食物链模型展示稳定、周期振荡与更复杂不规则的长期行为，并在问题四中用最大李雅普诺夫指数识别混沌候选区间；最后，将工业污染和外来入侵物种纳入情景扩展。针对问题一，本文进一步引入四大家鱼总量与单网次捕获量代理（CPUE*）作为可观测指标，其中 2025 年“四大家鱼恢复到禁渔前约 1.8 倍”作为硬校准锚点，湖北段“单网次捕获量提升 120\%”仅作为拟合外的一致性检验；结果表明，单区校准后四大家鱼总量在 2025 年达到基线的 1.805 倍，CPUE* 达到 1.838 倍，对应相对误差分别为 0.26\% 和 16.47\%，而模型中的水草末值降至 0.204 倍，提示在当前参数组下恢复伴随底层承压上升。基线修复情景下，江豚和长江鲟末期规模分别达到初值的 1.133 倍和 1.140 倍；在污染与入侵叠加情景下，内部综合健康指数由 0.686 降至 0.575，且权重扰动后前两类情景排序存在交换，提示该指标并非完全稳健。据此，本文建议将“禁渔”与“水质治理、通道修复、定向增殖放流和入侵物种压制”联动实施。
\end{abstract}

\newpage
\tableofcontents
\newpage
\pagenumbering{arabic}
\pagestyle{fancy}
